\clearpage
\begin{trapbox}
\textbf{\textcolor{CTrap}{易错点一｜只看等式，不看定义域}}
\begin{center}\begin{tikzpicture}[font=\small,>=stealth,line join=round]
\draw[gray!40,rounded corners=3pt] (0,0) rectangle (14.2,8.3);\node at (7.1,7.700000000000001) {缺少对应点：半条直线不关于原点对称};
\begin{scope}[shift={(4.1,3.9499999999999997)},x=1.0cm,y=0.7666666666666666cm]\draw[->,gray] (-3,0)--(3,0) node[right]{$x$};\draw[->,gray] (0,-3)--(0,3) node[above]{$y$};\draw[blue,thick,solid,domain=0:2.7,samples=180] plot(\x,{\x});\fill[orange!85!black] (0,0) circle(2.2pt);\fill[orange!85!black] (1,1) circle(2.2pt);\end{scope}\node[anchor=west,align=left,text width=5.0cm] at (8.4,4.0) {$f=x,\quad D=[0,+\infty)$\\[9pt]图上有 $(1,1)$\\[9pt]没有 $(-1,-1)$\\[9pt]只验证 $t=0$ 无法判定};
\node[text width=12.6cm,align=center] at (7.1,.65) {原点在图上也不够；必须让所有点转半圈后仍在原图像上。};
\end{tikzpicture}\end{center}



例如 $f(x)=x$ 只在 $[0,+\infty)$ 上有定义。若只对两边都存在的 $t$ 检查
$f(t)+f(-t)=0$，实际上只能检查 $t=0$；等式没有暴露问题。
但点 $(1,1)$ 的对应点 $(-1,-1)$ 不在图像上，所以图像\textbf{不关于原点中心对称}。
\textbf{改法：}先验证“有右边就必须有左边”，再检查等式。

\medskip

\end{trapbox}
\clearpage
\begin{trapbox}
\textbf{\textcolor{CTrap}{易错点二｜把轴对称与点对称混成一条式子}}
\begin{center}\begin{tikzpicture}[font=\small,>=stealth,line join=round]
\draw[gray!40,rounded corners=3pt] (0,0) rectangle (14.2,8.5);\node at (7.1,7.9) {同样取中间横坐标2：轴对称看同高，点对称看平均高度};
\node at (2.4,6.9) {轴：$f=(x-2)^2+3$};\begin{scope}[shift={(0.7,3.0500000000000003)},x=0.85cm,y=0.45cm]\draw[->,gray] (0,0)--(4,0) node[right]{$x$};\draw[->,gray] (0,-1)--(0,7) node[above]{$y$};\draw[blue,thick,solid,domain=0.5:3.5,samples=160] plot(\x,{(\x-2)^2+3});\fill[orange!85!black] (1,4) circle(2pt);\fill[orange!85!black] (3,4) circle(2pt);\end{scope}\node[align=center,text width=4cm] at (2.4,1.7) {$f(1)=f(3)=4$};\node at (7.050000000000001,6.9) {中心：$f=(x-2)^3+3$};\begin{scope}[shift={(5.3500000000000005,3.0500000000000003)},x=0.85cm,y=0.45cm]\draw[->,gray] (0,0)--(4,0) node[right]{$x$};\draw[->,gray] (0,-1)--(0,7) node[above]{$y$};\draw[violet,thick,solid,domain=0.5:3.5,samples=160] plot(\x,{(\x-2)^3+3});\fill[orange!85!black] (1,2) circle(2pt);\fill[orange!85!black] (3,4) circle(2pt);\end{scope}\node[align=center,text width=4cm] at (7.050000000000001,1.7) {$f(1)+f(3)=2+4=6$};\node at (11.7,6.9) {平移后：$g=t^3$};\begin{scope}[shift={(11.7,4.4)},x=0.85cm,y=0.45cm]\draw[->,gray] (-2,0)--(2,0) node[right]{$t$};\draw[->,gray] (0,-4)--(0,4) node[above]{$y$};\draw[orange!85!black,thick,solid,domain=-1.5:1.5,samples=160] plot(\x,{(\x)^3});\fill[orange!85!black] (-1,-1) circle(2pt);\fill[orange!85!black] (1,1) circle(2pt);\end{scope}\node[align=center,text width=4cm] at (11.7,1.7) {$g(-1)=-g(1)$};
\node[text width=12.6cm,align=center] at (7.1,.65) {前两图坐标范围相同；第三图是平移后的坐标范围，纵向显示 $[-4,4]$。};
\end{tikzpicture}\end{center}


\[
\begin{array}{ll}
\text{关于直线 }x=a\text{ 对称：}& f(a+t)=f(a-t),\\[3pt]
\text{关于点 }(a,b)\text{ 对称：}& f(a+t)+f(a-t)=2b.
\end{array}
\]
\textbf{改法：}轴对称是\textbf{左右一样高}；点对称是\textbf{左右高度以 $b$ 为中间数}。

\medskip

\end{trapbox}
\clearpage
\begin{trapbox}
\textbf{\textcolor{CTrap}{易错点三｜把 $2b$ 当成中心纵坐标}}
\begin{center}\begin{tikzpicture}[font=\small,>=stealth,line join=round]
\draw[gray!40,rounded corners=3pt] (0,0) rectangle (14.2,8.3);\node at (7.1,7.700000000000001) {左右点不必同高，但连线的中点始终是中心};
\begin{scope}[shift={(1.491304347826087,2.4166666666666665)},x=1.3043478260869565cm,y=0.5111111111111111cm]\draw[->,gray] (-0.3,0)--(4.3,0) node[right]{$x$};\draw[->,gray] (0,-1.5)--(0,7.5) node[above]{$y$};\draw[blue,thick,solid,domain=0.4:3.6,samples=180] plot(\x,{(\x-2)^3+3});\fill[orange!85!black] (1,2) circle(2.2pt);\fill[orange!85!black] (3,4) circle(2.2pt);\draw[orange!85!black] (1,2)--(3,4);\draw[gray,dashed] (2,-1.2)--(2,7.2);\draw[gray,dashed] (-.2,3)--(4.2,3);\fill[violet] (2,3) circle(2.6pt);\end{scope}\node[anchor=west,align=left,text width=5.0cm] at (8.4,4.0) {$f=(x-2)^3+3$\\[9pt]$P(3,4),\ Q(1,2)$\\[9pt]$\frac{3+1}{2}=2,\quad\frac{4+2}{2}=3$\\[9pt]中点 $C(2,3)$\\[9pt]$f(2+t)+f(2-t)=6$};
\node[text width=12.6cm,align=center] at (7.1,.65) {这条具体曲线示范一般公式；严格判定仍要求所有点都有对应点。};
\end{tikzpicture}\end{center}


$f(2+t)+f(2-t)=6$ 的中心是 $(2,3)$，\textbf{不是 $(2,6)$}。
因为 $6$ 是两点纵坐标的\textbf{和}，要除以 $2$。

\medskip

\end{trapbox}
\clearpage
\begin{trapbox}
\textbf{\textcolor{CTrap}{易错点四｜认为任何中心必须在图像上}}
\begin{center}\begin{tikzpicture}[font=\small,>=stealth,line join=round]
\draw[gray!40,rounded corners=3pt] (0,0) rectangle (14.2,8.3);\node at (7.1,7.700000000000001) {中心可以不在图上：两支曲线仍绕它成对};
\begin{scope}[shift={(2.1,2.4166666666666665)},x=1.0cm,y=0.5111111111111111cm]\draw[->,gray] (-1,0)--(5,0) node[right]{$x$};\draw[->,gray] (0,-1.5)--(0,7.5) node[above]{$y$};\draw[blue,thick,solid,domain=-0.5:1.7,samples=180] plot(\x,{1/(\x-2)+3});\draw[blue,thick,solid,domain=2.3:4.5,samples=180] plot(\x,{1/(\x-2)+3});\fill[orange!85!black] (1,2) circle(2.2pt);\fill[orange!85!black] (3,4) circle(2.2pt);\draw[gray,dashed] (2,-1.4)--(2,7.4);\draw[gray,dashed] (-.9,3)--(4.9,3);\draw[violet,fill=white] (2,3) circle(3pt);\end{scope}\node[anchor=west,align=left,text width=5.0cm] at (8.4,4.0) {$f=\frac1{x-2}+3$\\[9pt]中心 $(2,3)$\\[9pt]$D=\mathbb R\setminus\{2\}$\\[9pt]$f(1)=2,\quad f(3)=4$\\[9pt]但 $f(2)$ 不存在};
\node[text width=12.6cm,align=center] at (7.1,.65) {紫色空心定位点不属于函数图像；两条虚线是渐近线，不是曲线。};
\end{tikzpicture}\end{center}


$f(x)=1/(x-2)+3$ 的中心是 $(2,3)$，但 $x=2$ 不在定义域内。
\textbf{只有当 $a\in D$ 时，才可推出 $f(a)=b$。}

\medskip

\end{trapbox}
\clearpage
\begin{trapbox}
\textbf{\textcolor{CTrap}{易错点五｜只验证一两对数值就下结论}}
\begin{center}\begin{tikzpicture}[font=\small,>=stealth,line join=round]
\draw[gray!40,rounded corners=3pt] (0,0) rectangle (14.2,8.3);\node at (7.1,7.700000000000001) {某一对高度之和为6，不代表每一对都为6};
\begin{scope}[shift={(1.7000000000000002,1.9566666666666666)},x=1.2cm,y=0.6133333333333333cm]\draw[->,gray] (-0.5,0)--(4.5,0) node[right]{$x$};\draw[->,gray] (0,-0.5)--(0,7) node[above]{$y$};\draw[blue,thick,solid,domain=0:4,samples=180] plot(\x,{(\x-2)^2+2});\fill[orange!85!black] (0,6) circle(2.2pt);\fill[orange!85!black] (1,3) circle(2.2pt);\fill[orange!85!black] (3,3) circle(2.2pt);\fill[orange!85!black] (4,6) circle(2.2pt);\end{scope}\node[anchor=west,align=left,text width=5.0cm] at (8.4,4.0) {$f=(x-2)^2+2$\\[9pt]$f(1)+f(3)=3+3=6$\\[9pt]但 $f(0)+f(4)=6+6=12$\\[9pt]不关于 $(2,3)$ 中心对称};
\node[text width=12.6cm,align=center] at (7.1,.65) {这里只否定候选中心 $(2,3)$；一对点通过检验不等于全部配对通过。};
\end{tikzpicture}\end{center}


$f(1)+f(3)=6$ 只说明两点的平均高度是 $3$，并不能保证所有
$f(2+t)+f(2-t)$ 都等于 $6$。\textbf{完整判定必须对每一对都成立}。

\medskip

\end{trapbox}
\clearpage
\begin{trapbox}
\textbf{\textcolor{CTrap}{易错点六｜误以为点对称会改变单调方向}}
\begin{center}\begin{tikzpicture}[font=\small,>=stealth,line join=round]
\draw[gray!40,rounded corners=3pt] (0,0) rectangle (14.2,8.3);\node at (7.1,7.700000000000001) {横纵顺序同时反转，中心对称两侧同向单调};
\begin{scope}[shift={(1.2428571428571429,2.8)},x=1.4285714285714288cm,y=0.575cm]\draw[->,gray] (-0.1,0)--(4.1,0) node[right]{$x$};\draw[->,gray] (0,-2)--(0,6) node[above]{$y$};\draw[blue,thick,solid,domain=0.5:3.5,samples=180] plot(\x,{(\x-2)^3+2});\fill[orange!85!black] (0.5,-1.375) circle(2.2pt);\fill[orange!85!black] (1,1) circle(2.2pt);\fill[orange!85!black] (3,3) circle(2.2pt);\fill[orange!85!black] (3.5,5.375) circle(2.2pt);\end{scope}\node[anchor=west,align=left,text width=5.0cm] at (8.4,4.0) {示例 $f=(x-2)^3+2$\\[9pt]中心 $(2,2)$\\[9pt]左侧 $0.5<1$\\[9pt]右侧对应 $3.5>3$\\[9pt]$f(0.5)<f(1)$ 仍成立};
\node[text width=12.6cm,align=center] at (7.1,.65) {图取例6的一个具体函数；证明用比较式，适用于所有满足题设的函数。};
\end{tikzpicture}\end{center}


对于左右相互对应的两个区间，中心对称会\textbf{同时颠倒横纵顺序}，
所以右边递增，对应的左边也递增；右边递减，左边也递减。
注意只有在\textbf{相互对应的区间}比较时才能直接使用。

\medskip

\end{trapbox}
\clearpage
\begin{trapbox}
\textbf{\textcolor{CTrap}{易错点七｜把“有两个中心”都当成有周期}}
\begin{center}\begin{tikzpicture}[font=\small,>=stealth,line join=round]
\draw[gray!40,rounded corners=3pt] (0,0) rectangle (14.2,8.3);\node at (7.1,7.700000000000001) {两个中心不在同一高度：可以产生斜向平移};
\begin{scope}[shift={(3.6,3.5666666666666664)},x=1.0cm,y=0.7666666666666666cm]\draw[->,gray] (-2.5,0)--(3.5,0) node[right]{$x$};\draw[->,gray] (0,-2.5)--(0,3.5) node[above]{$y$};\draw[blue,thick,solid,domain=-2:3,samples=180] plot(\x,{\x});\fill[orange!85!black] (0,0) circle(2.2pt);\fill[orange!85!black] (1,1) circle(2.2pt);\end{scope}\node[anchor=west,align=left,text width=5.0cm] at (8.4,4.0) {$f(x)=x$\\[9pt]中心 $(0,0)$ 与 $(1,1)$\\[9pt]两次半转平移 $(2,2)$\\[9pt]纵坐标同时改变\\[9pt]不是水平周期};
\node[text width=12.6cm,align=center] at (7.1,.65) {直线有许多对称中心，但没有非零周期；不能遗漏两个中心同高的条件。};
\end{tikzpicture}\end{center}


若中心是 $(a,b)$ 和 $(c,b)$，同样高，可以推出周期 $2|c-a|$。
但 $y=x$ 上 $(0,0)$ 和 $(1,1)$ 都是对称中心，\textbf{它并不是周期函数}；
因为两个中心不在同一水平线上。

\medskip

\end{trapbox}
\clearpage
\begin{trapbox}
\textbf{\textcolor{CTrap}{易错点八｜混淆“一个中心”与“唯一中心”}}
\begin{center}\begin{tikzpicture}[font=\small,>=stealth,line join=round]
\draw[gray!40,rounded corners=3pt] (0,0) rectangle (14.2,8.3);\node at (7.1,7.700000000000001) {同一高度的两个中心，半转两次得到水平平移};
\begin{scope}[shift={(2.463636363636364,3.9499999999999997)},x=0.5454545454545454cm,y=1.5333333333333332cm]\draw[->,gray] (-2.5,0)--(8.5,0) node[right]{$x$};\draw[->,gray] (0,-1.5)--(0,1.5) node[above]{$y$};\draw[blue,thick,solid,domain=-2:8,samples=180] plot(\x,{sin(deg(pi*\x/3))});\fill[orange!85!black] (1,0.8660254037844386) circle(2.2pt);\fill[orange!85!black] (-1,-0.8660254037844386) circle(2.2pt);\fill[orange!85!black] (7,0.8660254037844386) circle(2.2pt);\end{scope}\node[anchor=west,align=left,text width=5.0cm] at (8.4,4.0) {图取 $f=\sin\frac{\pi x}{3}$\\[9pt]中心 $(0,0)$、$(3,0)$\\[9pt]$1\to-1\to7$\\[9pt]高度依次 $\frac{\sqrt3}2,-\frac{\sqrt3}2,\frac{\sqrt3}2$\\[9pt]相差 $6$，回到同高};
\node[text width=12.6cm,align=center] at (7.1,.65) {正文对任意函数的配对关系推出 $f(x+6)=f(x)$；6只保证是一个周期。};
\end{tikzpicture}\end{center}


公式能证明 $(a,b)$ \textbf{是一个}对称中心，未必是唯一的。
例如 $f(x)=\sin x$ 既关于 $(0,0)$ 又关于 $(\pi,0)$ 对称；
不能未经证明就声称“只有这个中心”。
\end{trapbox}
